% TOP_PROBLEM: {"id": 10000019, "title": "Isoperimetric dimension and nontrivial percolation threshold", "category_id": 19, "category": {"id": 19, "name": "probability", "display_name": "Probability", "description": "Problems involving probability theory, stochastic processes, and random structures.", "slug": "probability", "order_index": 19, "created_at": "2026-07-31T15:26:25.671Z"}, "status": "open", "problem_number": "AMR-099-0019", "set_id": 13, "statusQualification": "The source’s weaker infinite-dimensional case is already covered by the cited dimension-greater-than-four theorem; the full greater-than-one assertion is retained."}
% FIRST_POSTED: 2026-10-01
% ENTRY_KIND: statement_only
% UnsolvedMath ID 10000019; source catalogue code AMR-099-0019
% !TeX program = lualatex
% Definition and sources only; this file is not an LLM solution attempt.
\documentclass[11pt,a4paper]{article}
\usepackage[margin=25mm]{geometry}
\usepackage{fontspec}
\setmainfont{lmroman10-regular.otf}[BoldFont=lmroman10-bold.otf,
 ItalicFont=lmroman10-italic.otf,BoldItalicFont=lmroman10-bolditalic.otf]
\usepackage{amsmath,amssymb,mathtools}
\usepackage{unicode-math}
\setmathfont{latinmodern-math.otf}
\AtBeginDocument{\let\setminus\smallsetminus}
\usepackage{xurl,hyperref}
\hypersetup{colorlinks=true,urlcolor=blue}
\setcounter{secnumdepth}{0}
\setlength{\parindent}{0pt}
\setlength{\parskip}{5pt}
\setlength{\emergencystretch}{3em}
\begin{document}
\section*{Isoperimetric dimension and nontrivial percolation threshold}\label{10000019}
{\small UnsolvedMath ID 10000019 · AMR-099-0019 · Probability · AMR Open Problem Lists}
\subsection{Definitions and mathematical statement}
Let $G=(V,E)$ be an infinite connected graph with uniformly bounded degree. For a finite vertex set $A$, let $\partial_G A$ be its external vertex boundary. Its isoperimetric dimension is
\[
\operatorname{Dim}(G)=\sup\left\{s\ge1:\exists c>0\ \forall A\subset V,
\ 0<|A|<\infty,\quad |\partial_G A|\ge c|A|^{(s-1)/s}\right\}.
\]
For Bernoulli bond percolation with parameter $p$, independently retain each edge with probability $p$. If $o$ is any fixed vertex, define
\[
p_c(G)=\inf\{p\in[0,1]:\mathbb P_p(o\text{ lies in an infinite open component})>0\}.
\]
Connectedness makes this threshold independent of the root. Prove the implication $\operatorname{Dim}(G)>1\Longrightarrow p_c(G)<1$.

The isoperimetric inequality is required for every finite set, with a constant independent of its position and cardinality. No transitivity, polynomial growth, planarity, or symmetry is assumed. The desired parameter below one may depend on the entire graph; no universal upper bound for all such graphs is requested.

The source separately proposes the weaker case $\operatorname{Dim}(G)=\infty$. That case is included in the established theorem of Duminil-Copin, Goswami, Raoufi, Severo and Yadin: dimension greater than four suffices. Thus the unresolved target retained here is the full implication above one, in particular the range $1<\operatorname{Dim}(G)\le4$. This distinction separates the full target from the already established weaker high-dimensional case.
\subsection{Short English statement}
Does any uniform isoperimetric dimension strictly above one guarantee percolation after deleting a sufficiently small positive fraction of edges?
\subsection{Sources}
\begin{itemize}
\item[S1] ULAM.AI and the UnsolvedMath Contributors, supplied problems.json, record 10000019 (AMR-099-0019), AMR Open Problem Lists. Catalogue identity and source transcription; CC BY 4.0.
\url{https://huggingface.co/datasets/ulamai/UnsolvedMath}.
\item[S2] Benjamini - Coarse Geometry and Randomness (Saint-Flour notes), Conjecture 4.14, PDF page 34. Original source indexed by AMR-099-0019.
\url{https://www.wisdom.weizmann.ac.il/~itai/stflouraug24.pdf#page=34}.
\item[S3] I. Benjamini, Coarse Geometry and Randomness, primary Saint-Flour notes; the numbered question and its surrounding definitions.
\url{https://arquivo.pt/noFrame/replay/20201231041548id_/http://www.wisdom.weizmann.ac.il/~itai/stflouraug24.pdf}.
\item[S4] H. Duminil-Copin et al., Existence of phase transition for percolation using the Gaussian Free Field, Theorem 1.2 (dimension greater than four).
\url{https://arxiv.org/abs/1806.07733}.
\end{itemize}
\end{document}
